\documentclass[10 pt]{article}
\usepackage[utf8]{inputenc}
\usepackage{parskip}
\usepackage[utf8]{inputenc}
\usepackage[spanish]{babel}
\usepackage[left=2cm, right=2cm, top=1.5cm, bottom=2cm]{geometry}
\usepackage{float}
\usepackage{graphicx}
\usepackage{caption}
\usepackage{cancel}
\usepackage{amsmath,amsthm,amsfonts,amscd,amssymb,amsbsy}

\title{\textbf{\underline{Los enteros}}}
\author{}
\date{}

\begin{document}

\pagestyle{empty}
\maketitle
\thispagestyle{empty}

\section*{}

\LARGE{\textbf{\underline{Escritura e identificación}}}\\ \\

\Large{Hasta ahora hemos trabajado con los números enteros mediante representantes de sus respectivas clases de equivalencia. La suma, la multiplicación y la relación de orden se definieron directamente sobre dichos representantes, y comprobamos que ninguna de estas definiciones depende de los representantes escogidos. Esto sugiere la posibilidad de elegir, para cada entero, un representante ``distinguido'' o ``canónico'':} \\

\begin{itemize}
    \item \Large{\textbf{Teorema:} Sea $[(a,b)] \in \mathbb Z$. Si $[(a,b)] \neq [(0,0)]$, entonces existe $n \in \mathbb N \setminus \{0\}$ tal que $[(a,b)]=[(n,0)]$ o $[(a,b)]=[(0,n)]$.} \\
\end{itemize}

\Large{\textbf{Demostración:} Como $[(a,b)] \neq [(0,0)]$, resulta que $a+0 \neq 0+b$. De modo pues que $a \neq b$. Luego, por tricotomía, $a<b$ o $b<a$. Consideremos ambos casos:} \\

\Large{($\ast$) Supongamos que $a<b$. Entonces, como se demostró anteriormente, existe $p \in \mathbb N \setminus \{0\}$ tal que $b=a+p$. Con lo cual:}
\begin{align*}
    b&=b \\
    a+p&=b+0 \\
    &=0+b
\end{align*}
\Large{Por lo tanto, $[(a,b)]=[(0,p)]$.} \\

\Large{($\ast$) Similarmente: supongamos que $b<a$. Entonces, como se demostró anteriormente, existe $p \in \mathbb N \setminus \{0\}$ tal que $a=b+p$. Con lo cual:}
\begin{align*}
    a&=a \\
    a+0&=b+p \\
    &=p+b
\end{align*}
\Large{Por lo tanto, $[(a,b)]=[(p,0)]$.} \\
\begin{flushright}
   $\square$
\end{flushright}

\Large{El resultado anterior nos acerca a la representación habitual de los números enteros. El siguiente paso consiste en explicar cómo los números naturales se identifican con un subconjunto de $\mathbb Z$, como cabe esperar de una construcción cuyo propósito es extender el conjunto de los números naturales.} \\

\begin{itemize}
    \item \Large{\textbf{Definición:} La siguiente función se conoce como \textbf{inyección canónica} de $\mathbb N$ en $\mathbb Z$:} 
    \begin{align*}
        i:\mathbb N &\rightarrow \mathbb Z \\
        n &\mapsto [(n,0)]
    \end{align*}
\end{itemize}
\begin{flushright}
    $\blacksquare$
\end{flushright}

\Large{El papel de la inyección canónica es identificar de manera natural a $\mathbb N$ con un subconjunto de $\mathbb Z$. Para comprobar que esta identificación es adecuada, verificaremos que la aplicación es inyectiva y que preserva la suma, la multiplicación y la relación de orden:} \\

\begin{itemize}
    \item \Large{\textbf{Proposición:} La inyección canónica de $\mathbb N$ en $\mathbb Z$ es inyectiva.} \\
\end{itemize}

\Large{\textbf{Demostración:} Sean $n,m \in \mathbb N$, y supongamos que $i(n)=i(m)$. Es decir que $[(n,0)]=[(m,0)]$, y por tanto que $n+0=0+m$. Con lo cual, $n=m$. Tenemos entonces que $i$ es inyectiva.} \\

\begin{flushright}
    $\square$
\end{flushright}

\Large{La proposición anterior garantiza que la identificación de $\mathbb N$ en $\mathbb Z$ no colapsa información: \textit{cada número natural queda representado por un único número entero}. Para que esta identificación preserve la estructura de $\mathbb N$ debemos comprobar que la inyección canónica preserva las operaciones y la relación de orden. Comenzaremos verificando que preserva la suma y la multiplicación:} \\

\begin{itemize}
    \item \Large{\textbf{Proposición:} Sean $n,m \in \mathbb N$. \\

    i) $i(n+m)=i(n)+i(m)$. \\

    ii) $i(n \cdot m)=i(n) \cdot i(m)$.} \\
\end{itemize}

\Large{\textbf{Demostración:} \\

i) Observe que:}
\begin{align*}
    i(n+m)&=[(n+m,0)] \\
    &=[(n,0)]+[(m,0)] \\
    &=i(n)+i(m) 
\end{align*}
\Large{Observe que:}
\begin{align*}
    i(n \cdot m)&=[(n \cdot m,0)] \\
    &=[(n \cdot m,0 \cdot 0,n \cdot0+0 \cdot m)] \\
    &=[(n,0)] \cdot [(m,0)] \\
    &=i(n) \cdot i(m)
\end{align*}

\begin{flushright}
    $\square$
\end{flushright}

\Large{Veamos ahora que la inyección canónica también preserva el orden:} \\

\begin{itemize}
    \item \Large{\textbf{Proposición:} Sean $n,m \in \mathbb N$. \\

    i) $n<m$, si y sólo si, $i(n)<i(m)$. \\

    ii) $n \leq m$, si y sólo si, $i(n) \leq i(m)$.} \\
\end{itemize}

\Large{\textbf{Demostración:} \\

i) ($\rightarrow$) Supongamos que $n<m$. Entonces $n+0<m+0$, y por tanto $[(n,0)]<[(m,0)]$. Es decir que $i(n)<i(m)$. \\

($\leftarrow$) Supongamos que $i(n)<i(m)$. Es decir que $[(n,0)]<[(m,0)]$, y por tanto que $n+0<m+0$. Luego $n<m$. \\

ii) ($\rightarrow$) Supongamos que $n \leq m$. De modo pues que $n=m$ o $n<m$. Si $n=m$, claramente $i(n)=i(m)$. Y si $n<m$, entonces $i(n)<i(m)$ (por la parte anterior). Como podemos ver, en cualquier caso se sigue que $i(n) \leq i(m)$. \\

($\leftarrow$) Supongamos que $i(n) \leq i(m)$. De modo pues que $i(n)=i(m)$ o $i(n)<i(m)$. Si $i(n)=i(m)$, entonces $n=m$ (porque $i$ es inyectiva). Y si $i(n)<i(m)$, entonces $n<m$ (por la parte anterior). Como podemos ver, en cualquier caso se sigue que $n \leq m$. } \\

\begin{flushright}
    $\square$
\end{flushright}

\Large{Las proposiciones anteriores justifican que podamos identificar a $\mathbb N$ con un subconjunto de $\mathbb Z$. Introducimos a continuación una notación para dicha identificación:} \\

\begin{itemize}
    \item \Large{\textbf{Definición:} $\mathbb N_1:=i(\mathbb N)$. Es decir que: \\

    -- $\mathbb N_1=\{x \in \mathbb Z \hspace{0.2 cm}|\hspace{0.2 cm}(\exists n \in \mathbb N)(x=[(n,0)])\}$.} \\
\end{itemize}

\begin{flushright}
    $\blacksquare$
\end{flushright}

\Large{La preservación de las operaciones permite trasladar inmediatamente a $\mathbb N_1$ cualquier propiedad algebraica de $\mathbb N$. Veamos un ejemplo:} \\

\begin{itemize}
    \item \Large{\textbf{Proposición:} Sean $n,m \in \mathbb N$. Si $i(n)+i(m)=[(0,0)]$, entonces $i(n)=i(m)=[(0,0)]$.} \\
\end{itemize}

\Large{\textbf{Demostración:} Note que $i(0)=[(0,0)]$. Con lo cual:}
\begin{align*}
    i(n)+i(m)&=[(0,0)] \\
    i(n+m)&=i(0)
\end{align*}
\Large{Como $i$ es inyectiva, se sigue que $n+m=0$. Luego, como se demostró anteriormente, $n=m=0$. Por lo tanto, $i(n)=i(m)=i(0)=[(0,0)]$.} \\

\begin{flushright}
    $\square$
\end{flushright}

\Large{Este fenómeno no es exclusivo de la propiedad anterior. En general, cualquier propiedad algebraica de los números naturales expresada en términos de la suma, la multiplicación y la igualdad puede trasladarse a $\mathbb N_1$ mediante la inyección canónica. Del mismo modo, las demostraciones por inducción pueden realizarse indistintamente en $\mathbb N$ o en $\mathbb N_1$: demostrar un resultado en uno de estos conjuntos equivale a demostrar el correspondiente en el otro:} \\

\begin{itemize}
    \item \Large{\textbf{Teorema:} Sea $A \subseteq \mathbb N$. \\

    i) Si $A$ es inductivo, entonces $i(A)=\mathbb N_1$. \\

    ii) Si $[(0,0)] \in i(A)$ y $[(n,0)]+[(1,0)] \in i(A)$, para todo $[(n,0)] \in i(A)$, entonces $A=\mathbb N$ y $i(A)=\mathbb N_1$.} \\
\end{itemize}

\Large{\textbf{Demostración:} \\

i) De acuerdo con la hipótesis, $A=\mathbb N$ (por el principio de inducción). Luego $i(A)=i(\mathbb N)=\mathbb N_1$. \\

ii) Por hipótesis, $i(0) \in i(A)$ (porque $i(0)=[(0,0)]$). Como $i$ es inyectiva, necesariamente $0 \in A$. A continuación, consideremos $n \in A$. Como $[(n,0)] \in i(A)$ (porque $i(n) \in i(A)$), resulta que $[(n,0)]+[(1,0)] \in i(A)$. Es decir que $[(n+1,0)] \in i(A)$, y por tanto que $i(n+1) \in i(A)$. Y como $i$ es inyectiva, concluimos que $n+1 \in A$. Tenemos entonces que $A$ es inductivo, y por tanto que $A=\mathbb N$. Luego, por la parte anterior, $i(A)=\mathbb N_1$.} \\

\begin{flushright}
    $\square$
\end{flushright}

\Large{Algo completamente análogo ocurre con el orden:} \\

\begin{itemize}
    \item \Large{\textbf{Proposición:} Sea $A \subseteq \mathbb N_1$. Si $A \neq \emptyset$, entonces existe $a \in \mathbb N_1$ tal que $a=\min(A)$.} \\
\end{itemize}

\Large{En otras palabras: \textit{al análogo del principio del buen orden se cumple en $\mathbb N_1$}.} \\

\Large{\textbf{Demostración:} Consideremos el siguiente conjunto: \\

-- $A^{-1}=\{n \in \mathbb N \hspace{0.2 cm}|\hspace{0.2 cm}i(n) \in A\}$. \\

Como $A$ es no vacío, existe $x \in A$. Pero $A \subseteq \mathbb N_1$, así que $x \in \mathbb N_1$. Con lo cual, existe $n \in \mathbb N$ tal que $x=i(n)$. Tenemos entonces que $i(n) \in A$, y por tanto que $n \in A^{-1}$. De modo pues que $A^{-1} \neq \emptyset$. Luego, por el principio del buen orden, existe $m \in \mathbb N$ tal que $m=\min(A^{-1})$. A continuación, consideremos $i(k) \in A$, para algún $k \in \mathbb N$. Por definición, $k \in A^{-1}$. Luego $m \leq k$, y por tanto $i(m) \leq i(k)$. Por último, recordemos que $m \in A^{-1}$, así que $i(m) \in A$ (por definición). Por lo tanto, $i(m)=\min(A)$.} \\

\begin{flushright}
    $\square$
\end{flushright}

\Large{Desde el punto de vista algebraico, de orden y lógico, no existe diferencia entre trabajar en $\mathbb N$ y hacerlo en $\mathbb N_1$: trabajaremos indistintamente en uno u otro, según convenga. En consecuencia, a partir de este momento identificaremos ambos conjuntos y consideraremos a $\mathbb N$ como un subconjunto de $\mathbb Z$. Por ello, escribiremos ``$\mathbb N$'' en lugar de ``$\mathbb N_1$''. \\

La construcción de los números enteros ha cumplido ya su propósito: nos ha permitido establecer rigurosamente sus operaciones y su relación de orden. Además, la existencia de representantes canónicos muestra que todo número entero puede representarse de manera única mediante un número natural o el inverso aditivo de uno. En consecuencia, utilizaremos esta representación como la forma habitual de escribir los números enteros y recurriremos a las clases de equivalencia únicamente cuando resulte necesario.} \\

\Large{La \textbf{resta} se define como es usual en la teoría de anillos:} \\

\begin{itemize}
    \item \Large{\textbf{Definición:} Sean $n,m \in \mathbb Z$. Entonces $n-m:=n+(-m)$. Nos referimos a $n-m$ como la \textbf{resta} de $n$ y $m$.} \\
\end{itemize}

\Large{\textbf{Observación:} Se sigue de inmediato de la definición que $n-m \in \mathbb Z$.} \\

\begin{flushright}
    $\blacksquare$
\end{flushright}

\Large{Con esta notación, cada entero puede interpretarse como una diferencia de números naturales:} \\

\begin{itemize}
    \item \Large{\textbf{Proposición:} Sean $[(a,b)] \in \mathbb Z$. Entonces $[(a,b)]=i(a)-i(b)$.} \\
\end{itemize}

\Large{\textbf{Demostración:} Por definición:}
\begin{align*}
    [(a,b)]&=[(a,0)]+[(0,b)] \\
    &=[(a,b)]+(-[(b,0)]) \\
    &=i(a)+[-i(b)] \\
    &=i(a)-i(b)
\end{align*}

\begin{flushright}
    $\square$
\end{flushright}

\Large{Gracias a la identificación de $\mathbb N$ con un subconjunto de $\mathbb Z$, toda pareja de números naturales puede restarse sin restricciones. El resultado pertenece siempre a $\mathbb Z$, aunque puede no pertenecer a $\mathbb N$. El último resultado de la sección anterior nos permite determinar exactamente cuándo la diferencia de dos números enteros es un número natural:} \\

\begin{itemize}
    \item \Large{\textbf{Proposición:} Sean $n,m \in \mathbb Z$. Entonces $n-m \in \mathbb N$, si y sólo si, $m \leq n$.} \\
\end{itemize}

\begin{flushright}
    $\blacksquare$
\end{flushright}

\begin{itemize}
    \item \Large{\textbf{Corolario:} Sea $n \in \mathbb Z$. Entonces $n \in \mathbb N$, si y sólo si, $0 \leq n$.} \\
\end{itemize}

\begin{flushright}
    $\blacksquare$
\end{flushright}

\Large{Ahora podemos introducir las nociones de \textbf{entero positivo} y \textbf{entero negativo} de la manera esperada:} \\

\begin{itemize}
    \item \Large{\textbf{Definición:} Sea $\mathbb N^+:=\mathbb N \setminus \{0\}$. Adicionalmente, consideremos $n \in \mathbb Z \setminus \{0\}$. \\

    -- Diremos que $n$ es un \textbf{entero positivo}, si y sólo si, $n \in \mathbb N^+$. \\

    -- Diremos que $n$ es un \textbf{entero negativo}, si y sólo si, $n \notin \mathbb N^+$.} \\
\end{itemize}

\Large{\textbf{Observación:} De acuerdo con la definición, el cero no es positivo ni negativo.} \\

\begin{flushright}
    $\blacksquare$
\end{flushright}

\Large{\textbf{Ejemplos:}} \\

\Large{Sean $n,m \in \mathbb Z$.} \\

\begin{itemize}
    \item \Large{Supongamos que $[(n,m)]$ es un entero negativo. Tenemos entonces que $n<m$. Luego, como se demostró anteriormente, existe $p \in \mathbb N^+$ tal que $[(n,m)]=[(0,p)]$. En la práctica, nos referimos a $[(n,m)]$ simplemente como ``$-p$''.} \\

    \item \Large{Si $n=m$, entonces $[(n,m)]=[(0,0)]$. Nos referimos a $[(n,m)]$ simplemente como ``$0$''.} \\ 

    \item \Large{Supongamos que $[(n,m)]$ es un entero positivo. Tenemos entonces que $m<n$. Luego, como se demostró anteriormente, existe $q \in \mathbb N^+$ tal que $[(n,m)]=[(q,0)]$. En la práctica, nos referimos a $[(n,m)]$ simplemente como ``$q$''.} \\
\end{itemize}

\begin{flushright}
    $\blacksquare$
\end{flushright}

\Large{Para representar los números enteros podemos seguir utilizando el sistema de numeración posicional de los números naturales; únicamente es necesario anteponer el signo menos cuando el entero sea negativo.} \\

\Large{\textbf{Ejemplos:}} \\

\begin{itemize}
    \item \Large{Nos referimos $[(8,9)]$ como ``$-1$''.} \\

    \item \Large{Nos referimos a $[(100,200)]$ como ``$-100$''.} \\

    \item \Large{Nos referimos $[(-1,-1)]$ como ``$0$''.} \\

    \item \Large{Nos referimos a $[(5,4)]$ como ``$1$''.} \\

    \item \Large{Nos referimos a $[(50,25)]$ como ``$25$''.} \\
\end{itemize}

\begin{flushright}
    $\blacksquare$
\end{flushright}

\Large{Con esto concluye nuestro estudio de los sistemas numéricos. En adelante supondremos conocidas las propiedades fundamentales de los números \textbf{racionales} y \textbf{reales}, pues no será necesario desarrollar sus construcciones. A partir de este momento centraremos nuestra atención en las propiedades aritméticas de los números enteros, dando así inicio al estudio de la \textbf{teoría elemental de número}s.} \\

\end{document}